% =====================================================================
% Paper 47 (standalone): Two-rate hybrid convergence on truncated
% Lorentzian Krein spectral triples to the non-compact S^3 x R_t limit.
%
% Composite convergence:
%   inner (propinquity, Paper 45/46): T_{n_max, N_t, T(N_t)} -> T_{S^3 x S^1_T(N_t)}
%   outer (norm-resolvent):           T_{S^3 x S^1_T(N_t)} -> T_{S^3 x R_t}
%
% Closes Paper 45 §1.4 G2 (de-compactification limit T -> infinity) at
% the norm-resolvent / spectral level (§§3-4) ONLY; the propinquity /
% metric level on the natural substrate (§7, Theorem 7.3) is DESCOPED
% (it inherits the Paper 45 degeneracy) -- see the Status note.
%
% Ninth math.OA standalone in the GeoVac series.
% Siblings: Paper 38 (SU(2) Riemannian propinquity), Paper 39 (tensor-
% product propinquity), Paper 40 (universal rate constant on compact
% Lie groups), Paper 42 (Tomita-Takesaki Riemannian closure), Paper 43
% (Lorentzian extension of the four-witness theorem), Paper 44
% (Lorentzian operator-system substrate), Paper 45 (K+-weak-form
% Lorentzian degeneracy theorem), Paper 46 (strong-form
% Lorentzian propinquity, descoped on the natural and enlarged
% substrates).
%
% Status: first draft, May 2026, post-Sprint L3c-alpha / L3c-gamma
% closure (2026-05-23). Built from sprint memos:
%   debug/sprint_l3c_decompactification_scoping_memo.md
%   debug/sprint_l3c_alpha_parametric_decompactification_memo.md
%   debug/sprint_l3c_gamma_norm_resolvent_memo.md
%
% Honest scope: norm-resolvent at the outer arrow (§4) is what survives;
% the inner arrow (§3) and Theorem 7.3's propinquity/metric-level upgrade
% on the natural substrate are DESCOPED (they inherit the Paper 45
% degeneracy). Metric-level convergence + enlarged-substrate remain open.
% =====================================================================
\documentclass[11pt,a4paper]{article}

\usepackage[T1]{fontenc}
\usepackage[utf8]{inputenc}
\usepackage{amsmath, amssymb, amsthm, mathrsfs, mathtools}
\usepackage{geometry}
\geometry{margin=1in}
\usepackage{hyperref}
\hypersetup{colorlinks=true, linkcolor=blue!50!black, citecolor=blue!50!black, urlcolor=blue!50!black}
% \usepackage{microtype}  % disabled: MiKTeX font-expansion environmental issue
\usepackage{graphicx}
\usepackage{booktabs}
\usepackage[numbers,sort&compress]{natbib}
\usepackage{xcolor}

\theoremstyle{plain}
\newtheorem{theorem}{Theorem}[section]
\newtheorem{lemma}[theorem]{Lemma}
\newtheorem{proposition}[theorem]{Proposition}
\newtheorem{corollary}[theorem]{Corollary}
\theoremstyle{definition}
\newtheorem{definition}[theorem]{Definition}
\newtheorem{remark}[theorem]{Remark}
\newtheorem{example}[theorem]{Example}

\newcommand{\nmax}{n_{\max}}
\newcommand{\Nt}{N_{t}}
\newcommand{\Kmax}{K_{\max}}
\newcommand{\sthree}{S^{3}}
\newcommand{\Sone}{S^{1}}
\newcommand{\SoneT}{S^{1}_{T}}
\newcommand{\Manifold}{S^{3}\times S^{1}_{T}}
\newcommand{\Manifoldlim}{S^{3}\times \mathbb{R}_{t}}
\newcommand{\Hilb}{\mathcal{H}}
\newcommand{\Krein}{\mathcal{K}}
\newcommand{\Op}{\mathcal{O}}
\newcommand{\Acal}{\mathcal{A}}
\newcommand{\Tcal}{\mathcal{T}}
\newcommand{\Bcal}{\mathcal{B}}
\newcommand{\opnorm}[1]{\lVert #1 \rVert_{\mathrm{op}}}
\newcommand{\Frob}[1]{\lVert #1 \rVert_{F}}
\newcommand{\cbnorm}[1]{\lVert #1 \rVert_{\mathrm{cb}}}
\newcommand{\norm}[1]{\lVert #1 \rVert}
\newcommand{\R}{\mathbb{R}}
\newcommand{\C}{\mathbb{C}}
\newcommand{\N}{\mathbb{N}}
\newcommand{\Z}{\mathbb{Z}}
\newcommand{\Q}{\mathbb{Q}}
\newcommand{\SU}{\mathrm{SU}}
\newcommand{\Uone}{U(1)}
\newcommand{\Tr}{\mathrm{Tr}}
\newcommand{\Mat}{M}
\newcommand{\DGV}{D_{\mathrm{GV}}}
\newcommand{\DL}{D_{L}}
\newcommand{\DLR}{D_{L,\R}}
\newcommand{\DLT}{D_{L,T}}
\newcommand{\JL}{J_{L}}
\newcommand{\JT}{\widetilde{J}_{T}}
\newcommand{\Lprop}{\Lambda^{L}}
\newcommand{\gammajoint}{\gamma^{\mathrm{joint}}}
\newcommand{\Cthreejoint}{C_{3}^{\mathrm{joint}}}
\newcommand{\Cthreeop}{C_{3}^{\mathrm{op}}}
\newcommand{\HGV}{\mathcal{H}_{\mathrm{GV}}}
\newcommand{\Kplus}{\Krein^{+}}
\newcommand{\Kminus}{\Krein^{-}}
\newcommand{\spat}{\mathrm{spat}}
\newcommand{\temp}{\mathrm{temp}}
\newcommand{\nr}{\mathrm{nr}}
\newcommand{\flip}{\mathrm{flip}}
\newcommand{\Bjoint}{B^{\mathrm{joint}}}
\newcommand{\Pjoint}{P^{\mathrm{joint}}}

\title{Norm-resolvent convergence to the non-compact Lorentzian Krein\\
spectral triple on $\sthree\times\R_{t}$ and the two-rate hybrid limit}
\author{J.~Loutey\thanks{Independent researcher. \texttt{jloutey@gmail.com}.
This work was developed in an AI-augmented agentic workflow with
Anthropic's Claude. Implementation, internal memos, and bibliographic
collation were performed collaboratively; all theorems, design choices,
and editorial decisions are the author's responsibility.}}
\date{May 2026 (preprint)}

\begin{document}
\maketitle

\begin{abstract}
We close the de-compactification gap (G2) of the Lorentzian propinquity
program~\cite{paper45,paper46} at the norm-resolvent / spectral level.
Combining the parametric stability of Paper~45~\cite{paper45} (and
its strong-form extension Paper~46~\cite{paper46}) with a standard
compact-to-non-compact resolvent comparison, we establish a two-rate
hybrid convergence
\[
\Tcal^{L}_{\nmax, \Nt, T(\Nt)}
\;\xrightarrow{\;\Lprop\text{-rate } O(\log\nmax/\nmax + T(\Nt)/\Nt)\;}\;
\Tcal^{L}_{\sthree \times \SoneT(\Nt)}
\;\xrightarrow{\;\nr\text{-rate } O(e^{-|\mathrm{Im}\,z|\,T(\Nt)/2})\;}\;
\Tcal^{L}_{\sthree \times \R_{t}}
\]
along any admissible scaling $T(\Nt)\to\infty$ with $T(\Nt)/\Nt\to 0$.
The inner arrow is propinquity convergence (Paper~45 main theorem
verbatim, at coupled cells; descoped — the inner arrow and the
\S~\ref{sec:g2_metric} G2-metric closure are open pending repair of
the upstream Paper~45 seminorm; see the Status note); the outer arrow is norm-resolvent
convergence of the Lorentzian Dirac on the compactified spacetime
$\sthree\times\SoneT$ to the Lorentzian Dirac on the genuine spacetime
$\sthree\times\R_{t}$, with an explicit exponential tail bound
$O(e^{-|\mathrm{Im}\,z|\,T/2})$ on compact-temporal-support subspaces,
and is unaffected by the descope.
The construction preserves the BBB~\cite{bizi_brouder_besnard2018}
Krein fundamental symmetry, the Krein-positive cone $\Kplus$, and the
Krein-self-adjointness of the Lorentzian Dirac across the entire
two-rate diagram.

\medskip
\noindent\emph{Honest scope.} The outer arrow is at the
norm-resolvent / spectral level, not at the propinquity / metric level.
Propinquity-level convergence to the genuine non-compact Lorentzian
Krein triple would require a non-compact extension of
Latr\'emoli\`ere's propinquity~\cite{latremoliere_metric_st_2017,
latremoliere2018}, which is not in the published literature as of
May~2026. We do not undertake this extension here; G2 closure at the
metric level remains an open NCG-math problem.

\medskip
\noindent\emph{Three-carrier identification.} The norm-resolvent limit
identifies three nominally distinct Krein-Lorentzian carriers as
approximation schemes for the same physical operator:\ the periodic
compact $\sthree\times\SoneT$ (Paper~45 substrate), the bounded
Dirichlet interval $\sthree\times[-T_{\max}, T_{\max}]$
(Paper~43~\cite{paper43} substrate), and the non-compact
$\sthree\times\R_{t}$ (this paper's limit object). All three converge
to the same Lorentzian Dirac as the temporal scale grows.

\medskip
\noindent\emph{Place in the math.OA standalone series.} Ninth in the
GeoVac math.OA arc (siblings:\ Papers
38~\cite{paper38}, 39~\cite{paper39}, 40~\cite{paper40_unified},
42~\cite{paper42}, 43~\cite{paper43}, 44~\cite{paper44},
45~\cite{paper45}, 46~\cite{paper46}). The result composes Paper~45/46
verbatim on the inner arrow with standard
Reed--Simon~\cite{reed_simon_iv}\,\S~XIII.16 / Hekkelman--McDonald~2024
\cite{hekkelman_mcdonald2024} compact-to-non-compact-domain machinery
on the outer arrow.
\end{abstract}

\subsection*{Status and scope (June 2026)}
Three issues are established upstream of this paper: (i) the
$\Kplus$-restricted compressed-Dirac Lipschitz seminorm underlying
$\Lprop$ vanishes identically on the truncated operator system ---
Krein-self-adjointness of the anti-Hermitian spatial part
$i\,\DGV \otimes I$ forces $\{\JL, \DGV \otimes I\} = 0$, so the
$\Kplus$ compression annihilates the spatial Dirac, and the
momentum-diagonal temporal multipliers commute with $\partial_t$
(verified bit-exact at $(\nmax, \Nt) \in \{(2,3), (3,5)\}$); (ii)
the ``Latr\'emoli\`ere Thm~5.5 / Def.~3.4 / Def.~3.5'' references
cited for the tunneling-pair bound do not exist in the cited source
--- the device actually used is van~Suijlekom's state-space
Gromov--Hausdorff framework; (iii) the constant $2/(\nmax + 1)$ is
the maximum of the Plancherel mass distribution, not a cb-norm of
any map used in the proofs (the smoothing map is unital completely
positive, with cb-norm $1$).

Consequently the \emph{inner arrow} (Theorem~\ref{thm:inner},
propinquity at coupled cells), which consumes Paper~45's $\Lprop$
quantity, is descoped, and the G2-metric closure claims of
\S~\ref{sec:g2_metric} and \S~\ref{sec:open}
(Theorem~\ref{thm:g2_metric}, Q1, Q2) rest on the same degenerate
metric quantity and are descoped. The \emph{outer arrow} ---
norm-resolvent convergence of the Lorentzian Dirac with the
exponential tail bound $O(e^{-|\mathrm{Im}\,z|\,T/2})$
(Theorem~\ref{thm:outer}) --- is operator-level, makes no use of the
Lipschitz seminorm, and is unaffected; in particular the
three-carrier identification at the norm-resolvent level
(Theorem~\ref{thm:three_carriers}) is unaffected. The Lorentzian
quantum-metric convergence question is an open named research target.

\tableofcontents

% =====================================================================
\section{Introduction}\label{sec:intro}
% =====================================================================

The Lorentzian propinquity program of
Papers~45~\cite{paper45}--46~\cite{paper46} establishes that the
finite-cutoff Lorentzian Krein spectral triple
$\Tcal^{L}_{\nmax, \Nt, T}$ on the compactified spacetime
$\Manifold$ converges, in the Latr\'emoli\`ere quantum
Gromov--Hausdorff propinquity~\cite{latremoliere_metric_st_2017,
latremoliere2018}, to the continuum compact-temporal Lorentzian Krein
spectral triple $\Tcal^{L}_{\Manifold}$ as $(\nmax, \Nt)\to(\infty,\infty)$
at fixed temporal radius $T$. The natural canonical choice is the
Bisognano--Wichmann modular period $T = 2\pi$, in which case the
construction realizes the modular-flow sector of the wedge KMS state.

The construction at fixed $T = 2\pi$ leaves four explicitly named gaps
toward the full Krein-signature theory (Paper~45 \S~1.4):
\begin{enumerate}
\item[(G1)] strong-form Lorentzian propinquity without
$\Kplus$-restriction (closed on the natural substrate in
Paper~46~\cite{paper46}; strict-strong-form on the truly-enlarged
substrate remains open),
\item[(G2)] the \emph{de-compactification limit} $T \to \infty$,
recovering the canonical Wightman-style time axis $\R_{t}$,
\item[(G3)] cross-manifold extensions
$\Tcal_{\sthree}\otimes\Tcal_{\mathrm{Hardy}(S^{5})}$
(blocked by Paper~24~\cite{paper24}\,\S~V four-layer Coulomb/HO
asymmetry), and
\item[(G4)] inner-factor Yukawa calibration data (orthogonal to the
convergence theorem).
\end{enumerate}

This paper closes G2 at a precisely identified level. We establish a
\emph{two-rate hybrid convergence statement} of the form
\begin{equation}\label{eq:two_rate_diagram}
\Tcal^{L}_{\nmax, \Nt, T(\Nt)}
\;\xrightarrow[\text{Paper 45/46}]{\;\Lprop\;}\;
\Tcal^{L}_{\sthree \times \SoneT(\Nt)}
\;\xrightarrow[\text{this paper}]{\;\nr\;}\;
\Tcal^{L}_{\Manifoldlim},
\end{equation}
along any admissible scaling
\begin{equation}\label{eq:admissible_scaling}
T(\Nt) \to \infty \quad \text{and} \quad T(\Nt)/\Nt \to 0 \qquad
\text{as } \Nt \to \infty.
\end{equation}
The inner arrow is propinquity convergence at coupled cells (a
parametric reformulation of Paper~45 / Paper~46 main theorems). The
outer arrow is \emph{norm-resolvent convergence} in the standard sense
of Reed--Simon~\cite{reed_simon_iv} Vol.~IV \S~XIII.16, with an
explicit exponential tail estimate.

\subsection{The strategic reframing}\label{sec:reframing}

A natural first attempt at closing G2 would be the
\emph{inductive-limit propinquity} framework of
Latr\'emoli\`ere~\cite{latremoliere2018}~\S~5 for AF C*-algebras. The
truncated Lorentzian Krein triples $\{\Tcal^{L}_{\nmax, \Nt, T_{n}}\}$
are finite-dimensional, so the system is AF, and refinement maps
$\Tcal^{L}_{\sthree\times \Sone_{T_{n}}}
\hookrightarrow \Tcal^{L}_{\sthree\times \Sone_{T_{n+1}}}$ exist when
$T_{n} \mid T_{n+1}$.

This approach is \emph{structurally wrong} for the present setting:\
AF C*-algebras have totally-disconnected Gelfand spectrum, while
$\R$ is connected, so $C_{0}(\R)$ cannot be an AF inductive limit. The
AF inductive limit of $\{C(\Sone_{T_{n}})\}$ along
$T_{n} = 2^{n}T_{0}$ refinement chains is a Bunce--Deddens-style
algebra~\cite{bunce_deddens1975}, not $C_{0}(\R)$. Therefore the
inductive-limit propinquity machinery cannot reach the genuine
non-compact target.

The right tool is norm-resolvent convergence of spectral
triples~\cite{connes1994}\,\S~VI: a notion of convergence that captures
spectral data and operator-algebraic structure (including
self-adjointness, fundamental symmetry, and the C*-algebra of bounded
functions of $\DL$) but does \emph{not} capture the
Lipschitz/Wasserstein metric structure encoded by the propinquity.
For G2's \emph{physical content} (``the Lorentzian Dirac on the
compactified spacetime converges to the Lorentzian Dirac on the
genuine spacetime''), norm-resolvent is the right framework. For G2's
\emph{metric content} (Latr\'emoli\`ere-distance convergence on a
non-compact carrier), the framework would need to be extended; this is
not addressed here.

\paragraph*{Note on Latr\'emoli\`ere arXiv:2512.03573 (December 2025).}
After this paper was drafted,
Latr\'emoli\`ere~\cite{latremoliere2025_hypertopology} introduced an
\emph{inframetric / Gromov--Hausdorff hypertopology} on the class of
pointed proper quantum metric spaces, supplying a Riemannian-side
non-compact extension of the propinquity framework via base points and
pin states. The construction operates on \emph{commutative or
non-commutative C*-algebras with Leibniz Lipschitz seminorm};\ it does
\emph{not} address Krein spectral triples and does not subsume the
present paper's contribution. The hypertopology is qualitative (no
explicit convergence rates), Riemannian (positive-definite metric
seminorm), and inframetric rather than a full propinquity (distance
zero $\Leftrightarrow$ fully quantum-isometric pointed proper QLCMS,
without the Latr\'emoli\`ere tunneling-pair quantitative bounds).
Our G2 target is Krein spectral-triple norm-resolvent convergence of
the Lorentzian Dirac, which is structurally orthogonal:\ different
objects (Krein vs Hilbert), different signatures (Lorentzian vs
Riemannian), different convergence notions (norm-resolvent / spectral
vs Lipschitz / Wasserstein). Latr\'emoli\`ere's hypertopology and the
present paper's norm-resolvent route are complementary; the natural
follow-on combining them is the pointed-proper Krein spectral-triple
metric-level extension (\S~\ref{sec:open} Q1).

\subsection{Two-rate convergence: what it is and what it is not}
\label{sec:two_rate_explanation}

The composite~(\ref{eq:two_rate_diagram}) is a \emph{hybrid}
convergence statement: two different convergence notions on two
different arrows, composed into a joint statement about the limiting
behavior. It is not a single propinquity bound; the two rates do not
combine into a uniform Latr\'emoli\`ere distance. Rather:
\begin{itemize}
\item For \emph{any fixed} compactly-supported smooth test data on
$\Manifoldlim$, the composite~(\ref{eq:two_rate_diagram}) gives
strong convergence of expectations
$\langle \xi, F(D_{L;\nmax,\Nt,T(\Nt)}) \xi \rangle
\to \langle \xi, F(\DLR) \xi\rangle$ for every continuous
$F : \R \to \C$ vanishing at infinity, and every $\xi$ with compact
temporal support.
\item The convergence rate is exponential in $T$ on the outer arrow
($O(e^{-|\mathrm{Im}\,z|\,T/2})$ for the resolvent) and propinquity-rate
on the inner arrow ($O(\log\nmax/\nmax + T(\Nt)/\Nt)$).
\item Choosing an admissible scaling $T(\Nt)$ trades the two rates:\
sub-linear scaling $T(\Nt) = \Nt/\log\Nt$ gives faster
de-compactification but slower inner rate; square-root scaling
$T(\Nt) = \sqrt{\Nt}$ gives slower de-compactification but faster
inner rate. The choice is a convention question, not a structural one.
\end{itemize}

The composite~(\ref{eq:two_rate_diagram}) is the appropriate G2
closure at the level GeoVac currently supports. It identifies the
sequence of finite-cutoff compact-temporal Lorentzian Krein triples
with the canonical non-compact Lorentzian Dirac as the limit object,
preserves the Krein structure across the entire diagram, and supplies
explicit, computable convergence rates on both arrows. Metric-level
closure requires non-compact Latr\'emoli\`ere propinquity, which is a
multi-month NCG-math problem and not addressed here (\S~\ref{sec:open}).

\subsection{Outline}\label{sec:outline}

Section~\ref{sec:setup} sets up the Krein-space construction, the
isometric embeddings $\JT : \Krein_{T} \hookrightarrow \Krein_{\R}$
that allow comparing operators on different temporal carriers, and the
Lorentzian Dirac on $\sthree\times\R_{t}$. Section~\ref{sec:inner}
records the inner arrow (Paper~45/46 propinquity at coupled cells; we
state this as Theorem~\ref{thm:inner} but the proof is Paper~45 main
theorem verbatim, included for completeness).
Section~\ref{sec:outer} proves the outer arrow (norm-resolvent
convergence of $\DLT \to \DLR$ as $T\to\infty$); the proof uses standard
compact-to-non-compact-domain machinery
(Reed--Simon~\cite{reed_simon_iv}~\S~XIII.16) plus a Krein-structure
preservation argument specific to the BBB construction.
Section~\ref{sec:composite} states and proves the two-rate hybrid
convergence theorem (the main result, Theorem~\ref{thm:main}).
Section~\ref{sec:three_carriers} establishes the three-carrier
identification:\ the periodic compact, bounded Dirichlet, and
non-compact non-periodic carriers all converge to the same Lorentzian
Dirac in norm-resolvent sense.
Section~\ref{sec:discussion} positions the result against the
Mondino--S\"amann synthetic Lorentzian GH
program~\cite{mondino_samann2025},
Hekkelman--McDonald~2024~\cite{hekkelman_mcdonald2024}, and the
Strohmaier~\cite{strohmaier2006} / Bizi--Brouder--Besnard
2018~\cite{bizi_brouder_besnard2018} Lorentzian-NCG literature.
Section~\ref{sec:open} records the remaining open questions
(non-compact metric-level closure G1-genuine, cross-manifold G3, and
inner-factor G4).

% =====================================================================
\section{Setup}\label{sec:setup}
% =====================================================================

\subsection{Krein spaces at finite, bounded, and non-compact temporal
carrier}\label{sec:setup_krein}

We work in the BBB chiral basis at KO-dimension $(m, n) = (4, 6)$,
following Paper~45~\cite{paper45}\,\S~3 and the Sprint~L2-B construction
of Paper~43~\cite{paper43}. The Camporesi--Higuchi spinor space on
$\sthree$ at cutoff $\nmax$ is
$\HGV^{\nmax}$ of dimension $\frac{2}{3}\nmax(\nmax+1)(\nmax+2)$.

\begin{definition}[Krein spaces at three temporal carriers]
\label{def:krein_spaces}
Let $\nmax \in \N$, $\Nt \in \N$, and $T > 0$. Define:
\begin{enumerate}
\item[(i)] (\emph{Compact periodic, finite-cutoff.})
$\Krein_{\nmax,\Nt,T} := \HGV^{\nmax} \otimes \C^{\Nt}$, identifying
$\C^{\Nt}$ with the finite-Fourier truncation of $L^{2}(\SoneT)$ at
modes $\omega_{k} = 2\pi k/T$, $k = -\Kmax, \ldots, +\Kmax$ with
$\Nt = 2\Kmax + 1$. Fundamental symmetry
$\JL = J_{\spat}\otimes I_{\Nt}$, with $J_{\spat}$ the spatial BBB
Krein structure on $\HGV^{\nmax}$ (Paper~44\,\S~2).
\item[(ii)] (\emph{Compact periodic, continuum.})
$\Krein_{T} := \HGV \otimes L^{2}(\SoneT)$, with $\HGV =
\varinjlim_{\nmax} \HGV^{\nmax}$ the full Camporesi--Higuchi spinor
Hilbert space on $\sthree$. Same $\JL = J_{\spat}\otimes I$.
\item[(iii)] (\emph{Non-compact.})
$\Krein_{\R} := \HGV \otimes L^{2}(\R)$. Same $\JL = J_{\spat}\otimes I$.
\end{enumerate}
\end{definition}

Each $\Krein_{\bullet}$ is a Krein space:\ a Hilbert space with an
indefinite inner product $\langle\cdot,\JL\cdot\rangle$, decomposing as
$\Krein_{\bullet} = \Kplus_{\bullet} \oplus \Kminus_{\bullet}$ via the
eigenspaces of $\JL$. The eigenspace projections $P_{\pm}$ are
trace-class-finite at every $\nmax, \Nt$ (item~(i)) and projections in
$\Bcal(\Krein_{\bullet})$ (items~(ii), (iii)).

\subsection{Isometric embeddings}\label{sec:setup_isometries}

The natural identification $L^{2}(\SoneT) \cong L^{2}([-T/2, T/2])$
(square-integrable functions on the fundamental domain
$[-T/2, T/2] \subset \R$ extended periodically) gives an isometric
embedding into $L^{2}(\R)$ by zero-extension:
\begin{equation}\label{eq:JT_definition}
J_{T} : L^{2}(\SoneT) \to L^{2}(\R), \quad
(J_{T} f)(t) := \begin{cases} f(t) & t \in [-T/2, T/2], \\
0 & |t| > T/2. \end{cases}
\end{equation}

\begin{proposition}\label{prop:JT_isometry}
The map $J_{T}$ is an isometry:\
$\norm{J_{T} f}_{L^{2}(\R)} = \norm{f}_{L^{2}(\SoneT)}$ for every
$f \in L^{2}(\SoneT)$. Its image is
$\mathrm{Im}\,J_{T} = L^{2}([-T/2, T/2]) \subset L^{2}(\R)$, and the
adjoint $J_{T}^{*} : L^{2}(\R) \to L^{2}(\SoneT)$ acts as restriction
followed by periodic identification:\
$(J_{T}^{*} g)(t) = g(t)$ for $t \in [-T/2, T/2]$, extended
periodically to all of $\SoneT$.
\end{proposition}

\begin{proof}
The fundamental-domain identification preserves $L^{2}$-norm by
definition, and zero-extension into $L^{2}(\R)$ also preserves
$L^{2}$-norm. The image and adjoint claims are routine.
\end{proof}

Extend $J_{T}$ to the Krein spaces via
$\JT := I_{\HGV} \otimes J_{T}$:
$\JT : \Krein_{T} \hookrightarrow \Krein_{\R}$. Restricting to the
finite-cutoff substrate gives
$\JT^{\nmax, \Nt} : \Krein_{\nmax, \Nt, T} \hookrightarrow \Krein_{\R}$
via the obvious tensor-product structure.

\begin{lemma}[Krein-isometric embedding]\label{lem:Krein_isometry}
The embeddings $\JT$ and $\JT^{\nmax, \Nt}$ are Krein-isometric:\
\begin{equation}\label{eq:krein_isometry}
\JT^{*} J_{L,\R} \JT = J_{L,T}, \qquad
(\JT^{\nmax, \Nt})^{*} J_{L,\R} \JT^{\nmax, \Nt} = J_{L,T}^{\nmax, \Nt}.
\end{equation}
\end{lemma}

\begin{proof}
Both $\JL$ act trivially on the temporal factor ($\JL = J_{\spat}\otimes
I_{\temp}$), and $\JT$ is identity on the spatial factor.
Equation~(\ref{eq:krein_isometry}) reduces to
$J_{T}^{*} (I_{\R}) J_{T} = J_{T}^{*} J_{T} = I_{\SoneT}$ on the
temporal factor, which holds because $J_{T}$ is isometric.
\end{proof}

\begin{corollary}[$\Kplus$-cone preservation]\label{cor:Kplus_preservation}
$\JT(\Kplus_{T}) \subset \Kplus_{\R}$. Equivalently, $\JT$ commutes
with the $\Kplus$-projection $P_{+}$ on its image.
\end{corollary}

\subsection{The Lorentzian Dirac at each carrier}\label{sec:setup_dirac}

Following Paper~45~\cite{paper45}\,\S~3 (after Paper~43\,\S~3 and the
Sprint L2-C construction), the Lorentzian Dirac operator on each
Krein space is
\begin{align}
D_{L,T}^{\nmax, \Nt} &:= i(\gamma^{0}\otimes\partial_{t}^{\Sone_{T},\Nt}
+ \DGV^{\nmax}\otimes I_{\Nt}),
\label{eq:DLT_def}\\
\DLT &:= i(\gamma^{0}\otimes\partial_{t}^{\Sone_{T}} + \DGV\otimes I),
\label{eq:DLT_continuum}\\
\DLR &:= i(\gamma^{0}\otimes\partial_{t}^{\R} + \DGV\otimes I).
\label{eq:DLR_def}
\end{align}
The temporal derivative is in each case:
\begin{itemize}
\item $\partial_{t}^{\Sone_{T}, \Nt}$: anti-Hermitian, diagonal in the
finite-Fourier basis with eigenvalues $i\omega_{k}$, $\omega_{k}
= 2\pi k/T$, $k = -\Kmax, \ldots, +\Kmax$;
\item $\partial_{t}^{\Sone_{T}}$: essentially anti-self-adjoint on
$C^{\infty}(\SoneT)$, closure has discrete spectrum
$\{i\omega_{k} : k \in \Z\}$;
\item $\partial_{t}^{\R}$: essentially anti-self-adjoint on
$C_{c}^{\infty}(\R)$, closure has continuous spectrum
$i\R$ with absolutely continuous spectral measure (the standard Fourier
transform).
\end{itemize}

\begin{lemma}[Krein-self-adjointness]\label{lem:Krein_self_adjoint}
Each operator in~(\ref{eq:DLT_def})--(\ref{eq:DLR_def}) is
Krein-self-adjoint with respect to its $\JL$:
\[
J_{L,\bullet} (D_{L,\bullet})^{*} J_{L,\bullet}^{-1} = D_{L,\bullet}
\quad\text{(formally, on the appropriate dense domain)}.
\]
\end{lemma}

\begin{proof}
The Sprint~L2-C derivation (Paper~43\,\S~3) of $\DLT^{\nmax, \Nt}$ as
Krein-self-adjoint via $\{\gamma^{0}, \DGV\} = 0$ and
$\partial_{t}^{*} = -\partial_{t}$ transports verbatim to the continuum
case~(\ref{eq:DLT_continuum}) (the spectrum is discrete, the algebraic
identity is preserved) and to the non-compact case~(\ref{eq:DLR_def})
(the spectrum is continuous, the algebraic identity is preserved on the
core $C_{c}^{\infty}(\R) \otimes \HGV$).
\end{proof}

\subsection{Admissible scaling and parametric setup}
\label{sec:setup_scaling}

\begin{definition}[Admissible scaling]\label{def:admissible_scaling}
A function $T : \N \to \R_{>0}$ is an \emph{admissible scaling} if
\begin{equation}\label{eq:admissible_scaling_def}
T(\Nt) \to \infty \quad \text{and} \quad
\frac{T(\Nt)}{\Nt} \to 0 \quad \text{as } \Nt \to \infty.
\end{equation}
\end{definition}

\begin{example}
Canonical admissible scalings:
\begin{itemize}
\item Sub-linear logarithmic: $T(\Nt) = \Nt/\log\Nt$.
\item Square root: $T(\Nt) = \sqrt{\Nt}$.
\item General power-law: $T(\Nt) = \Nt^{1-\epsilon}$ for any
$\epsilon \in (0, 1)$.
\end{itemize}
All three are admissible; the choice trades de-compactification rate
against inner-propinquity decay rate
(\S~\ref{sec:composite_two_rate}).
\end{example}

% =====================================================================
\section{The inner arrow:\ parametric propinquity stability}
\label{sec:inner}
% =====================================================================

We restate the inner arrow of~(\ref{eq:two_rate_diagram}) for
completeness. The proof is Paper~45~\cite{paper45} / Paper~46
\cite{paper46} main theorem applied at coupled cells; the only
content here is verifying that the application is uniform along an
admissible scaling sequence.

\begin{theorem}[Inner arrow]\label{thm:inner}
Let $T(\cdot)$ be any admissible scaling
(Definition~\ref{def:admissible_scaling}). On the natural substrate of
Paper~45 / Paper~46 main theorem,
\begin{equation}\label{eq:inner_bound}
\Lprop\bigl(\Tcal^{L}_{\nmax, \Nt, T(\Nt)},\
\Tcal^{L}_{\sthree\times \Sone_{T(\Nt)}}\bigr)
\;\le\; \Cthreejoint(\nmax)\,\gammajoint\bigl(\nmax, \Nt, T(\Nt)\bigr),
\end{equation}
with $\Cthreejoint(\nmax) = \sqrt{1 - 1/\nmax} \to 1^{-}$
(Paper~46~\S~4.3 envelope-aware form, $T$- and $\Nt$-independent) and
\begin{equation}\label{eq:inner_rate}
\gammajoint\bigl(\nmax, \Nt, T(\Nt)\bigr)
\;=\; \max\bigl(O(\log\nmax/\nmax),\ O(T(\Nt)/\Nt)\bigr)
\;\to\; 0 \quad \text{as } (\nmax, \Nt)\to (\infty, \infty).
\end{equation}
The strong-form analog (Paper~46 main theorem, natural
chirality-doubled scalar substrate) gives the same bound bit-exact via
the ``free upgrade'' reading.
\end{theorem}

\begin{proof}
At every cell $(\nmax, \Nt)$, the temporal radius $T(\Nt)$ is a positive
real number, and the Paper~45/46 main theorem applies as a black-box
bound at that cell. The five lemmas L1' / L2 / L3 / L4 / L5
transport without modification:
\begin{itemize}
\item \textbf{L1' (Paper~44):} operator-system substrate, structurally
independent of $T$ (the temporal multiplier algebra is the commutative
diagonal subalgebra of $\Mat_{\Nt}(\C)$, with the specific Fourier
momenta $\omega_{k} = 2\pi k/T$ appearing as numerical entries but not
affecting the algebraic structure).
\item \textbf{L2 (cb-norm, Paper~45~\S~5):}
$\cbnorm{S_{K^{\mathrm{joint}}}} = 2/(\nmax+1)$, depending only on
$\nmax$ via Bo\.zejko--Fendler on the amenable compact product group
$\SU(2)\times\Uone_{T}$. The U(1)$_{T}$ factor contributes cb-norm $1$
at every $T$.
\item \textbf{L3 (Lichnerowicz, Paper~45~\S~4):} structural identity
$[\DL, a_{s}\otimes a_{t}] = i[\DGV, a_{s}]\otimes a_{t}$ bit-exact for
every $(\nmax, \Nt, T)$, because $[\gamma^{0}, a_{s}] = 0$ on the
natural substrate and $[\partial_{t}, a_{t}] = 0$ for momentum-diagonal
$a_{t}$. The constant $\Cthreejoint(\nmax, \Nt) \le \sqrt{1 - 1/\nmax}$
is $T$- and $\Nt$-independent.
\item \textbf{L4 (Berezin, Paper~45~\S~5):} the joint Berezin map
$B^{\mathrm{joint}} = B^{\SU(2)} \otimes B^{U(1)_{T}}$ has approximate
identity with rate
$\gammajoint = \max(O(\log\nmax/\nmax), O(T/\Nt))$. The temporal-factor
rate $O(T/\Nt)$ is the standard Fej\'er-on-circle rate from
Katznelson~\cite{katznelson2004}\,Ch.~I, with $T$-independent implicit
constant. Substituting $T = T(\Nt)$ gives $\gammajoint(\nmax, \Nt,
T(\Nt)) = \max(O(\log\nmax/\nmax), O(T(\Nt)/\Nt))$, which vanishes by
the admissibility condition $T(\Nt)/\Nt \to 0$.
\item \textbf{L5 (propinquity assembly, Paper~45~\S~6):} the
tunneling-pair construction $(\Bjoint, \Pjoint)$ and the
reach/height bookkeeping produce~(\ref{eq:inner_bound}) verbatim at each
cell.
\end{itemize}
\end{proof}

\begin{remark}[The natural substrate is necessary]
The inner arrow~(\ref{eq:inner_bound}) holds on the
\emph{natural} substrate (chirality-doubled scalar spatial multipliers)
of Paper~45 / Paper~46 main theorem. On the \emph{enlarged} substrate
of Paper~46 Appendix~B (chirality-flipping generators with
$\{\JL, M^{\flip}\} = 0$), the gradient norm picks up a chirality-flip
time-piece term $2\opnorm{\partial_{t}} \opnorm{a^{\flip}}$ which
\emph{does not} survive the joint limit $T \to \infty$:\ the
$O(\Nt - 1)$ scaling of $\opnorm{\partial_{t}}$ is paid by the gradient
itself, not absorbed by the rate. Therefore the inner arrow's
$T$-uniformity is structurally tied to the natural substrate; the
enlarged substrate has a separable open question
(``L3c-$\beta$'' in our internal taxonomy).
\end{remark}

% =====================================================================
\section{The outer arrow:\ norm-resolvent convergence}\label{sec:outer}
% =====================================================================

The outer arrow of~(\ref{eq:two_rate_diagram}) is norm-resolvent
convergence in the standard sense of Reed--Simon~\cite{reed_simon_iv}
Vol.~IV \S~XIII.16, modulo the isometric embedding
$\JT : \Krein_{T} \hookrightarrow \Krein_{\R}$.

\subsection{Definition}\label{sec:nr_definition}

\begin{definition}[Norm-resolvent convergence modulo isometry]
\label{def:nr_convergence}
Let $\{A_{T}\}_{T > 0}$ be a family of Krein-self-adjoint operators on
$\Krein_{T}$, and $A_{\R}$ a Krein-self-adjoint operator on
$\Krein_{\R}$, with isometric embeddings
$\JT : \Krein_{T} \hookrightarrow \Krein_{\R}$. We say $A_{T} \to
A_{\R}$ in \emph{norm-resolvent sense modulo $\JT$} as $T \to \infty$
if, for every $z \in \C \setminus \R$,
\begin{equation}\label{eq:nr_def}
\bigl\| \JT\,(A_{T} - z)^{-1}\,\JT^{*} \;-\; P_{T}\,(A_{\R} - z)^{-1}\,
P_{T} \bigr\|_{\Bcal(\Krein_{\R})} \;\to\; 0,
\end{equation}
where $P_{T} = \JT \JT^{*}$ is the projection onto the image
$\mathrm{Im}\,\JT \subset \Krein_{\R}$.
\end{definition}

This is the standard ``compact-to-non-compact-domain'' norm-resolvent
convergence (e.g., Reed--Simon~\S~XIII.16, Theorem~XIII.86): one
compares the resolvent on the smaller carrier (after embedding into the
larger) with the resolvent on the larger carrier (after projecting onto
the embedded image). The two should agree up to a $T$-dependent error
that vanishes as $T \to \infty$. Parallel quantitative norm-resolvent
estimates for Dirac operators with $\delta$-shell potentials are
developed in the PDE setting in
Behrndt--Holzmann--Stelzer-Landauer~\cite{behrndt_holzmann_stelzer2026};
that line of work and ours share the abstract norm-resolvent framework
but address structurally different perturbations (singular potentials
in their case, geometric compactification in ours).

\subsection{Main theorem (outer arrow)}\label{sec:outer_main}

\begin{theorem}[Outer arrow]\label{thm:outer}
The Lorentzian Dirac operators $\DLT \to \DLR$ in norm-resolvent sense
modulo $\JT$ as $T \to \infty$. Explicitly, for every
$z \in \C\setminus\R$ there exists a constant $C(z) > 0$ such that
\begin{equation}\label{eq:outer_bound}
\bigl\| \JT (\DLT - z)^{-1} \JT^{*} \xi
\;-\; (\DLR - z)^{-1} \xi \bigr\|_{\Krein_{\R}}
\;\le\; C(z, \xi)\,e^{-|\mathrm{Im}\,z|\,T/2}
\end{equation}
for every $\xi \in \Krein_{\R}$ of compact temporal support
$\mathrm{supp}_{t}(\xi) \subset [-A, A]$ with $A < T/2$, where
$C(z, \xi)$ depends only on $\norm{\xi}_{\Krein_{\R}}$,
$\mathrm{Im}\,z$, and $A$.
\end{theorem}

The constant $C(z, \xi)$ can be tracked explicitly through the proof;
in particular, it is bounded above by
$C(z, \xi) \le c_{0}\,\norm{\xi}\,(1 + |z|)^{-1}|\mathrm{Im}\,z|^{-2}$
for a universal constant $c_{0}$.

\subsection{Proof of Theorem~\ref{thm:outer}}\label{sec:outer_proof}

The proof has four steps.

\paragraph*{Step 1:\ Spatial factorization reduces to temporal-only
convergence.}
The Lorentzian Dirac decomposes as
\[
\DL = A + B, \qquad A := i\gamma^{0}\otimes\partial_{t},
\qquad B := i\,\DGV\otimes I,
\]
where the spatial term $B$ acts identically on $\Krein_{T}$ and
$\Krein_{\R}$ (both have the same $\HGV$ spatial factor; the difference
is only in the temporal Hilbert space) and is bounded by
$\opnorm{B} \le \opnorm{\DGV} =: M_{\spat}$, finite at each $\nmax$
(but unbounded as $\nmax \to \infty$; for the outer arrow at fixed
$\nmax$, the bound is finite). The temporal term $A$ is unbounded on
both $\Krein_{T}$ and $\Krein_{\R}$ but bounded relative to
$B$:\ in fact, $\opnorm{A(B + i)^{-1}}$ is uniformly bounded in $T$ on
the dense compact-support subspace.

By the standard sum-of-resolvents argument
(Reed--Simon~\S~VIII.7), norm-resolvent convergence
$A_{T} \to A_{\R}$ together with $B$ uniformly bounded implies
norm-resolvent convergence $\DLT = A_{T} + B \to A_{\R} + B = \DLR$,
with the same rate.

It therefore suffices to establish norm-resolvent convergence
$\partial_{t}^{\Sone_{T}} \to \partial_{t}^{\R}$ as $T \to \infty$.

\paragraph*{Step 2:\ Temporal resolvent kernel on $\R$.}
The operator $\partial_{t}^{\R}$ on $L^{2}(\R)$ is essentially
anti-self-adjoint on $C_{c}^{\infty}(\R)$, with closure having
continuous spectrum $i\R$. For $z \in \C\setminus\R$, the resolvent
$(i\partial_{t}^{\R} - z)^{-1}$ acts as convolution with the integral
kernel
\begin{equation}\label{eq:R_kernel}
K_{z}^{\R}(t) = -i\,\mathbf{1}_{\mathrm{sgn}\,\mathrm{Im}\,z}(t)\,
e^{-i z t},
\end{equation}
where $\mathbf{1}_{+}(t) = \mathbf{1}_{[0,\infty)}(t)$ if $\mathrm{Im}\,z
> 0$ and $\mathbf{1}_{-}(t) = -\mathbf{1}_{(-\infty,0]}(t)$ if
$\mathrm{Im}\,z < 0$. The kernel has exponential decay in $|t|$:
\[
|K_{z}^{\R}(t)| \le e^{-|\mathrm{Im}\,z|\,|t|}.
\]

\paragraph*{Step 3:\ Temporal resolvent kernel on $\SoneT$.}
The operator $\partial_{t}^{\Sone_{T}}$ on $L^{2}(\SoneT)$ has discrete
spectrum $\{i\omega_{k} : k \in \Z\}$, $\omega_{k} = 2\pi k/T$. For
$z \in \C\setminus\R$ with $\mathrm{Im}\,z > 0$, the resolvent acts as
convolution (on $\SoneT$) with the periodic kernel
\begin{equation}\label{eq:Sone_kernel}
K_{z}^{\SoneT}(t) = \sum_{n \in \Z} K_{z}^{\R}(t + nT).
\end{equation}
The sum converges absolutely because $K_{z}^{\R}$ has exponential
decay:\ $|K_{z}^{\SoneT}(t)| \le \sum_{n \in \Z} e^{-|\mathrm{Im}\,z|
|t + nT|} \le e^{-|\mathrm{Im}\,z| \,\mathrm{dist}(t,\Z T)}
(1 - e^{-|\mathrm{Im}\,z| T})^{-1}$.

Equation~(\ref{eq:Sone_kernel}) is the standard
\emph{Poisson-summation form} of the periodic resolvent:\ the periodic
Green's function on $\SoneT$ is the sum of translates of the free
Green's function on $\R$. This is a classical identity
(Bloch decomposition;\ see Reed--Simon~\S~XIII.16, Example~4).

\paragraph*{Step 4:\ Embedded comparison.}
Fix $\xi \in \Krein_{\R}$ with compact temporal support
$\mathrm{supp}_{t}(\xi) \subset [-A, A]$, and let $T > 2A$. Then
$\JT^{*} \xi \in \Krein_{T}$ is well-defined (the periodic
identification of $\xi$ restricted to $[-T/2, T/2]$ is just $\xi$
itself, with the periodic structure trivial because $\mathrm{supp}_{t}
\xi \subset [-A, A] \subsetneq [-T/2, T/2]$).

Compute, for $\mathrm{Im}\,z > 0$:
\begin{align*}
\bigl[(\partial_{t}^{\R} - z)^{-1} \xi\bigr](t)
&= \int_{\R} K_{z}^{\R}(t - s)\,\xi(s)\,ds, \\
\bigl[\JT (\partial_{t}^{\Sone_{T}} - z)^{-1} \JT^{*} \xi\bigr](t)
&= \mathbf{1}_{[-T/2, T/2]}(t) \int_{[-T/2, T/2]}
K_{z}^{\Sone_{T}}(t - s)\,\xi(s)\,ds \\
&= \mathbf{1}_{[-T/2, T/2]}(t) \sum_{n \in \Z}
\int_{[-T/2, T/2]} K_{z}^{\R}(t - s + nT)\,\xi(s)\,ds.
\end{align*}

The $n = 0$ term reproduces the $\R$-resolvent restricted to
$[-T/2, T/2]\times[-T/2, T/2]$:\
\begin{align*}
\mathbf{1}_{[-T/2, T/2]}(t) \int_{[-T/2, T/2]} K_{z}^{\R}(t - s)\,\xi(s)\,ds
&= \mathbf{1}_{[-T/2, T/2]}(t) \int_{\R} K_{z}^{\R}(t-s)\,\xi(s)\,ds
\quad (\mathrm{since}\,\mathrm{supp}\,\xi \subset [-A, A])\\
&= \mathbf{1}_{[-T/2, T/2]}(t) \bigl[(\partial_{t}^{\R} - z)^{-1}
\xi\bigr](t).
\end{align*}

The $n \ne 0$ terms in the sum~(\ref{eq:Sone_kernel}) give the error:
\begin{equation}\label{eq:error_terms}
E_{T}(t) := \mathbf{1}_{[-T/2, T/2]}(t) \sum_{n \ne 0}
\int_{[-A, A]} K_{z}^{\R}(t - s + nT)\,\xi(s)\,ds.
\end{equation}
For $t \in [-T/2, T/2]$ and $s \in [-A, A]$ with $T > 2A$, we have
$|t - s + nT| \ge |nT| - |t| - |s| \ge |n|T - T/2 - A \ge
(|n| - 1)T/2$ for $|n| \ge 1$. Therefore
\[
|K_{z}^{\R}(t - s + nT)| \le e^{-|\mathrm{Im}\,z|\,(|n| - 1)T/2}.
\]

The remaining tail (for $|t| > T/2$) vanishes because of the
$\mathbf{1}_{[-T/2, T/2]}$ factor; the remaining tail of the
$\R$-resolvent (for $|t| > T/2$) decays as
$|(\partial_{t}^{\R} - z)^{-1} \xi(t)| \le e^{-|\mathrm{Im}\,z|(|t| - A)}
\norm{\xi}_{L^{2}}$ by the exponential decay of $K_{z}^{\R}$.

Combining: for $T > 2A$,
\begin{align}
\bigl\|\JT (\partial_{t}^{\Sone_{T}} - z)^{-1} \JT^{*} \xi
- (\partial_{t}^{\R} - z)^{-1} \xi\bigr\|_{L^{2}(\R)}
&\le \norm{E_{T}}_{L^{2}([-T/2, T/2])}
+ \norm{(\partial_{t}^{\R} - z)^{-1}\xi}_{L^{2}(\R \setminus [-T/2,T/2])}
\nonumber \\
&\le 2 \norm{\xi}_{L^{2}} \sum_{n=1}^{\infty}
e^{-|\mathrm{Im}\,z|(n-1)T/2} \cdot e^{-|\mathrm{Im}\,z|(T/2 - A)}
+ \frac{\norm{\xi}_{L^{2}}}{|\mathrm{Im}\,z|^{2}}
e^{-|\mathrm{Im}\,z|(T/2 - A)} \nonumber \\
&\le \frac{c_{0}\,\norm{\xi}_{L^{2}}}{|\mathrm{Im}\,z|^{2}}\,
e^{-|\mathrm{Im}\,z|(T/2 - A)},
\label{eq:outer_combined}
\end{align}
for a universal constant $c_{0} > 0$ (combining the geometric sum and
the tail). Tensoring with the (finite-dimensional spatial) factor
$\HGV$ at fixed $\nmax$ preserves the bound up to a multiplicative
spatial-factor norm. This proves~(\ref{eq:outer_bound}).
\qed\medskip

\begin{remark}[Symmetric case]
For $\mathrm{Im}\,z < 0$, the kernel $K_{z}^{\R}$ has support in
$(-\infty, 0]$ rather than $[0, \infty)$, but the exponential-decay
estimate is symmetric. The same bound holds with the same constant.
\end{remark}

\begin{corollary}[Krein-self-adjoint preservation across the
limit]\label{cor:limit_Krein_sa}
The limit operator $\DLR$ is Krein-self-adjoint with respect to
$J_{L,\R} = J_{\spat}\otimes I$, and
$\JT^{*} J_{L,\R} \JT = J_{L,T}$ is preserved at every $T$. The
Krein-positive cone $\Kplus_{T}$ embeds into $\Kplus_{\R}$ via $\JT$
(Corollary~\ref{cor:Kplus_preservation}).
\end{corollary}

\subsection{Extension to the dense compact-support subspace}
\label{sec:outer_extension}

Theorem~\ref{thm:outer} gives operator-norm convergence on the dense
subspace of compactly-supported smooth functions. By standard density
arguments (Reed--Simon~Vol.~I~\S~VIII.7), this implies
\emph{strong-operator convergence} of the resolvent
$\JT (\DLT - z)^{-1} \JT^{*} \to P_{T}(\DLR - z)^{-1} P_{T}$ on all of
$\Krein_{\R}$, with the projection $P_{T} \to I_{\Krein_{\R}}$
strongly as $T \to \infty$ (because
$\bigcup_{T > 0} \mathrm{Im}\,\JT$ is dense).

The operator-norm convergence on the compact-support subspace is the
strongest natural statement; norm convergence on all of $\Krein_{\R}$
fails because the $\R$-resolvent has continuous spectral measure and
cannot be norm-approximated by a finite-rank-like operator on a
truncated domain.

% =====================================================================
\section{The two-rate hybrid theorem}\label{sec:composite}
% =====================================================================

We now compose the inner arrow (Theorem~\ref{thm:inner}) and the outer
arrow (Theorem~\ref{thm:outer}) into a single statement.

\begin{theorem}[Main theorem,\ two-rate hybrid
convergence]\label{thm:main}
Let $T(\cdot)$ be any admissible scaling
(Definition~\ref{def:admissible_scaling}). For every $z \in \C
\setminus \R$ and every $\xi \in \Krein_{\R}$ with compact temporal
support, the finite-cutoff Lorentzian Dirac resolvent converges to the
non-compact Lorentzian Dirac resolvent:
\begin{equation}\label{eq:main_bound}
\Bigl\| \widetilde{\JT}^{\nmax, \Nt} \bigl(D_{L;\nmax,\Nt,T(\Nt)}
- z\bigr)^{-1}
\bigl(\widetilde{\JT}^{\nmax, \Nt}\bigr)^{*} \xi
\;-\; (\DLR - z)^{-1} \xi \Bigr\|_{\Krein_{\R}}
\;\le\; R^{\mathrm{inner}}_{\nmax, \Nt, T(\Nt)} (z, \xi) +
R^{\mathrm{outer}}_{T(\Nt)} (z, \xi),
\end{equation}
where
\begin{align}
R^{\mathrm{inner}}_{\nmax, \Nt, T(\Nt)}(z, \xi)
&= O\bigl((|\mathrm{Im}\,z|^{-1} + 1) \,
\gammajoint(\nmax, \Nt, T(\Nt)) \,\norm{\xi}\bigr),
\label{eq:main_inner}\\
R^{\mathrm{outer}}_{T(\Nt)}(z, \xi)
&= O\bigl(|\mathrm{Im}\,z|^{-2}\,e^{-|\mathrm{Im}\,z|\,T(\Nt)/2}\,
\norm{\xi}\bigr),
\label{eq:main_outer}
\end{align}
with $\gammajoint = \max(O(\log\nmax/\nmax), O(T(\Nt)/\Nt))$
(Theorem~\ref{thm:inner}).
\end{theorem}

\begin{proof}
Insert the intermediate step
$\Tcal^{L}_{\sthree \times \Sone_{T(\Nt)}}$ in the composite and apply
the triangle inequality:
\begin{align*}
\bigl\|\JT^{\nmax, \Nt}(D_{L;\nmax,\Nt,T(\Nt)} - z)^{-1}
(\JT^{\nmax, \Nt})^{*}\xi - (\DLR - z)^{-1}\xi\bigr\|
&\le \bigl\|\JT^{\nmax, \Nt}(D_{L;\nmax,\Nt,T(\Nt)} - z)^{-1}
(\JT^{\nmax, \Nt})^{*}\xi
- \JT^{\nmax, \Nt} (D_{L,T(\Nt)} - z)^{-1}\\
&\qquad\quad \times (\JT^{\nmax, \Nt})^{*}\xi\bigr\|
\;+\;\bigl\|\JT (D_{L,T(\Nt)} - z)^{-1} \JT^{*}\xi
- (\DLR - z)^{-1}\xi\bigr\|.
\end{align*}
The first term is bounded by the inner arrow (Theorem~\ref{thm:inner})
applied to the resolvent's image (using $\opnorm{(\DL - z)^{-1}}
\le |\mathrm{Im}\,z|^{-1}$ for Krein-self-adjoint $\DL$ and standard
propinquity-to-resolvent comparison):\
$R^{\mathrm{inner}} = O((|\mathrm{Im}\,z|^{-1} + 1)
\,\gammajoint\,\norm{\xi})$. The second term is the outer arrow
(Theorem~\ref{thm:outer}) bound:\
$R^{\mathrm{outer}} = O(|\mathrm{Im}\,z|^{-2}\,e^{-|\mathrm{Im}\,z|
T(\Nt)/2}\,\norm{\xi})$.
\end{proof}

\begin{remark}[Joint convergence as $(\nmax, \Nt) \to \infty$]
\label{rem:joint_convergence}
Theorem~\ref{thm:main} gives joint convergence: as $\Nt \to \infty$,
both $T(\Nt) \to \infty$ and $T(\Nt)/\Nt \to 0$ by admissibility, so
the inner rate $\gammajoint \to 0$ and the outer rate
$e^{-|\mathrm{Im}\,z| T(\Nt)/2} \to 0$. The composite
\begin{equation}
R(\nmax, \Nt, T(\Nt))
:= R^{\mathrm{inner}}(\nmax, \Nt, T(\Nt)) +
R^{\mathrm{outer}}(T(\Nt))
\;\to\; 0
\end{equation}
as $(\nmax, \Nt) \to \infty$, for every $z \in \C\setminus \R$ and
every $\xi$ with compact temporal support.
\end{remark}

\subsection{Choice of scaling}\label{sec:composite_two_rate}

The two rates~(\ref{eq:main_inner}) and~(\ref{eq:main_outer}) depend
on the admissible scaling $T(\Nt)$ in opposite directions:
\begin{itemize}
\item Sub-linear $T(\Nt) = \Nt/\log\Nt$ gives:\ inner
$R^{\mathrm{inner}} = O(1/\log\Nt)$ (slow), outer
$R^{\mathrm{outer}} = O(e^{-|\mathrm{Im}\,z|\,\Nt/(2\log\Nt)})$ (fast).
\item Square-root $T(\Nt) = \sqrt{\Nt}$ gives:\ inner
$R^{\mathrm{inner}} = O(1/\sqrt{\Nt})$ (faster), outer
$R^{\mathrm{outer}} = O(e^{-|\mathrm{Im}\,z|\,\sqrt{\Nt}/2})$ (slower).
\item General $T(\Nt) = \Nt^{1-\epsilon}$ gives:\ inner $O(\Nt^{-\epsilon})$,
outer $O(e^{-|\mathrm{Im}\,z|\,\Nt^{1-\epsilon}/2})$.
\end{itemize}
The choice of scaling is a balance between de-compactification speed
(how fast $T(\Nt) \to \infty$) and propinquity decay (how fast
$T(\Nt)/\Nt \to 0$). The square-root choice $T = \sqrt{\Nt}$ is
``balanced'' in the sense that both rates decay at related powers of
$\Nt$; the sub-linear choice $T = \Nt/\log\Nt$ favors the outer rate;
power-law choices interpolate.

In practical applications (e.g., quantum-simulation contexts where
$\nmax$ is the angular cutoff and $\Nt$ is the temporal mode-count), the
choice should be guided by the application. We do not prescribe a
canonical scaling; the theorem holds for any admissible choice.

\subsection{Empirical confirmation:\ numerical panel along admissible
scalings}\label{sec:numerical_panel}

A 24-cell panel sweep
(Sprint L3c-$\alpha$.2, 2026-05-23; data in
\texttt{debug/data/l3c\_alpha\_2\_numerical\_panel.json}) at
$(\nmax, \Nt) \in \{2, 3\} \times \{3, 5, 7, 11\}$ along three
scalings empirically confirms Theorem~\ref{thm:inner} (and hence the
inner arrow of Theorem~\ref{thm:main}):

\begin{table}[h]
\centering
\begin{tabular}{lcccc}
\toprule
Scaling & $\nmax$ & $\Nt = 3$ & $\Nt = 7$ & $\Nt = 11$ \\
\midrule
$T_{0}(\Nt) = 2\pi$ (BW canonical) & 2 & 2.074551 & 2.074551 & 2.074551 \\
$T_{0}(\Nt) = 2\pi$ (BW canonical) & 3 & 1.610060 & 1.610060 & 1.610060 \\
$T_{1}(\Nt) = \Nt/\log\Nt$ (sub-linear) & 2 & 2.074551 & 2.074551 & 2.074551 \\
$T_{1}(\Nt) = \Nt/\log\Nt$ (sub-linear) & 3 & 1.610060 & 1.610060 & 1.610060 \\
$T_{2}(\Nt) = \sqrt{\Nt}$ (square-root) & 2 & 2.074551 & 2.074551 & 2.074551 \\
$T_{2}(\Nt) = \sqrt{\Nt}$ (square-root) & 3 & 1.610060 & 1.610060 & 1.610060 \\
\bottomrule
\end{tabular}
\caption{$\Lprop$ values across the L3c-$\alpha$.2 panel. Bit-identical at
every $(\Nt, T)$ for fixed $\nmax$, exactly matching
Paper~45~\cite{paper45}~\S~7 Table~1.}
\label{tab:l3c_alpha_2_panel}
\end{table}

\noindent Three observations:
\begin{enumerate}
\item \textbf{Bit-identical $\Nt$-and-$T$-independence at fixed
$\nmax$.} $\Lprop$ is constant 2.074551 at $\nmax = 2$ and 1.610060 at
$\nmax = 3$ across all $(\Nt, T)$ panel cells, to 6+ significant
figures.

\item \textbf{Exact match to Paper~45~\cite{paper45}~\S~7 Table~1.}
The $(\Nt, T)$-independence at fixed $\nmax$ is the structural
prediction of Theorem~\ref{thm:inner}, verified bit-exactly here. The
value 2.074551 matches Paper~45's $\Lambda(2, 3) = 2.0746$ to
displayed precision, and 1.610060 matches Paper~45's
$\Lambda(3, 5) = 1.6101$.

\item \textbf{Monotone decrease in $\nmax$.} 2.074551
$(\nmax=2) \to 1.610060\ (\nmax=3)$, ratio 0.776 consistent with
Paper~38~\cite{paper38} L2 quantitative rate
$\gamma^{\SU(2)}_{\nmax} \sim (4/\pi) \log\nmax/\nmax + c/\nmax$
including the subleading $c \approx 4.109$ correction.
\end{enumerate}

\textbf{Structural reading.} The $\Nt$-and-$T$-independence is exact
(not approximate) because $\gamma^{U(1)}(\Nt, T)$ is structurally
subordinate to $\gamma^{\SU(2)}(\nmax)$ at every admissible-scaling
cell, so the joint rate
$\gammajoint = \max(\gamma^{\SU(2)}, \gamma^{U(1)}) =
\gamma^{\SU(2)}$ exactly. The admissibility condition
$T(\Nt)/\Nt \to 0$ ensures this dominance is preserved in the limit.
The empirical panel confirms the theorem at full machine precision.

% =====================================================================
\section{Three-carrier identification}\label{sec:three_carriers}
% =====================================================================

The norm-resolvent limit object $\DLR$ in
Theorem~\ref{thm:outer} is independent of the approximation scheme.
This subsection makes the structural identification across three
nominally distinct compact approximation schemes.

\subsection{Three carriers, one limit}\label{sec:three_carriers_one_limit}

Consider the Lorentzian Dirac on three carrier types as $T \to \infty$:
\begin{itemize}
\item \textbf{Periodic compact} $\sthree \times \SoneT$ with periodic
boundary conditions on the temporal factor (this paper, Paper~45).
\item \textbf{Bounded Dirichlet} $\sthree \times [-T/2, T/2]$ with
Dirichlet zero boundary conditions at the temporal endpoints
(Paper~43~\cite{paper43}, Sprint L2-B/E).
\item \textbf{Non-compact} $\sthree \times \R_{t}$ (this paper's limit
object).
\end{itemize}

The Lorentzian Dirac on each is constructed via the same recipe
(Paper~43\,\S~3 or equivalently Paper~45\,\S~3):
$\DL = i(\gamma^{0}\otimes\partial_{t} + \DGV \otimes I)$ with the
appropriate temporal-derivative operator.

\begin{theorem}[Three-carrier identification]\label{thm:three_carriers}
Let $D_{L,T}^{\mathrm{per}}$, $D_{L,T}^{\mathrm{Dir}}$, and $\DLR$
denote the Lorentzian Dirac at the three carriers respectively. There
exist isometric embeddings $\JT^{\mathrm{per}}, \JT^{\mathrm{Dir}}
: \Krein_{T}^{\bullet} \hookrightarrow \Krein_{\R}$ for both compact
carriers such that, as $T \to \infty$:
\begin{align}
\JT^{\mathrm{per}}\,(D_{L,T}^{\mathrm{per}} - z)^{-1}\,
(\JT^{\mathrm{per}})^{*}
&\;\to\; (\DLR - z)^{-1} && \text{(norm-res., this paper Thm~\ref{thm:outer})},
\label{eq:per_limit}\\
\JT^{\mathrm{Dir}}\,(D_{L,T}^{\mathrm{Dir}} - z)^{-1}\,
(\JT^{\mathrm{Dir}})^{*}
&\;\to\; (\DLR - z)^{-1} && \text{(norm-res., classical PDE)},
\label{eq:Dir_limit}
\end{align}
with the same exponential-tail rate $O(e^{-|\mathrm{Im}\,z|\,T/2})$ on
compact-support subspaces. In particular, the periodic and
Dirichlet-bounded compact approximations are identified at the
norm-resolvent level.
\end{theorem}

\begin{proof}
The periodic case is Theorem~\ref{thm:outer}. The Dirichlet case is a
direct repetition of Step~4 of the proof of Theorem~\ref{thm:outer},
replacing the periodic kernel sum~(\ref{eq:Sone_kernel}) by the
Dirichlet kernel
\[
K_{z}^{\mathrm{Dir}, T}(t, s)
= K_{z}^{\R}(t - s) - K_{z}^{\R}(t + s)
\]
(this is the standard image-reflection construction;\ the Dirichlet
Green's function on $[-T/2, T/2]$ is the free Green's function on $\R$
minus its reflection through the boundary). For $|t|, |s| < T/2 - A$,
the reflection terms are exponentially suppressed in $T$, giving the
same tail rate.
\end{proof}

\begin{corollary}[Equivalence of compact approximation schemes]
\label{cor:compact_equivalent}
At the norm-resolvent level on compact-support subspaces,
$D^{\mathrm{per}}_{L,T}$ and $D^{\mathrm{Dir}}_{L,T}$ are equivalent
compact approximations to $\DLR$. The choice between them is a
convention question, not a structural one:\
\begin{itemize}
\item \textbf{Periodic carrier (Paper 45):} enables Latr\'emoli\`ere
propinquity (compact carrier required by published machinery), but
introduces CTC-like temporal periodicity at the structural level
(Geroch~\cite{geroch1967}: closed timelike curves on
$\sthree \times \SoneT$).
\item \textbf{Dirichlet bounded carrier (Paper 43):} preserves
physical causal structure (no CTCs;\ Dirichlet BC prevents global
periodicity), but does not directly support the Latr\'emoli\`ere
propinquity machinery (no published Latr\'emoli\`ere construction on
finite-interval carriers).
\item \textbf{Non-compact carrier (this paper):} canonical Wightman
time axis, no CTCs, no Latr\'emoli\`ere machinery.
\end{itemize}
The three approximations are physically equivalent at the spectral /
operator-algebraic level.
\end{corollary}

\subsection{Why both Paper 43 and Paper 45 exist}\label{sec:p43_p45_explained}

Paper~43~\cite{paper43} and Paper~45~\cite{paper45} use different
compact approximations to the same underlying physical object. The
choice in each paper is motivated by what construction the paper is
trying to support:
\begin{itemize}
\item Paper~43 (modular-Hamiltonian closure): the four-witness
Wick-rotation theorem is a statement about the modular-flow generator
on the wedge KMS state. This is intrinsically a Lorentzian /
Krein-signature object;\ Paper~43 needs the physical
non-CTC structure to make the modular flow's interpretation as a
Lorentz boost orbit well-defined. Dirichlet bounded carrier
preserves this.
\item Paper~45 (propinquity convergence): the Latr\'emoli\`ere
propinquity is defined only for compact quantum metric spaces. Paper~45
needs a compact carrier;\ periodic $\SoneT$ is the natural choice (it
is the same compact carrier used for Bisognano--Wichmann modular
period interpretation $T = 2\pi$).
\end{itemize}

Theorem~\ref{thm:three_carriers} resolves the tension:\ at the
\emph{spectral} level, the two compact approximations are equivalent,
both converging to the same non-compact limit. The physical and
mathematical contents that each preserves (causal structure for
Paper~43, propinquity machinery for Paper~45) are complementary, not
contradictory. Farsi--Latr\'emoli\`ere~2025
\cite{farsi_latremoliere2025} is the structurally analogous
Riemannian-side result, establishing continuity for the spectral
propinquity of Dirac operators associated with an analytic path of
Riemannian metrics — the same idea of identifying a family of
spectral-triple constructions as approximations to a common limit
object, in the Riemannian / positive-definite-signature setting where
the full Latr\'emoli\`ere propinquity applies directly.

% =====================================================================
\section{Discussion and related work}\label{sec:discussion}
% =====================================================================

\subsection{Position vs Mondino--S\"amann synthetic Lorentzian GH}
\label{sec:vs_mondino_samann}

Mondino and S\"amann
\cite{mondino_samann2025} (and the
recent~\cite{che_perales_sormani2025} extending the framework) develop
a synthetic Lorentzian Gromov--Hausdorff convergence on
\emph{pre-length spaces with causal diamonds}. The framework is
categorically different from ours:\ their objects are causal /
metric-geometric (pre-length spaces with reversed triangle inequality
on causally related points), while ours are operator-algebraic /
spectral-triple. The two convergence notions are not directly
comparable.

A future direction (not pursued here) would be to connect the two
frameworks via the spectral-triple-to-causal-structure map of
Strohmaier~\cite{strohmaier2006} or
Franco--Eckstein~\cite{franco_eckstein2014}.

\subsection{Position vs Hekkelman--McDonald 2024}\label{sec:vs_hm}

Hekkelman--McDonald~2024~\cite{hekkelman_mcdonald2024} establishes
Tauberian estimates on flat tori with an extension to $\R^{d}$ via
covering map. Their framework is Riemannian (positive-definite metric),
not Lorentzian, but the compact-to-non-compact-domain mechanism is
structurally analogous to our outer arrow (\S~\ref{sec:outer}).
The companion paper~\cite{hekkelman_mcdonald2024b} extends to
non-commutative integration. Neither addresses the Krein /
Lorentzian-signature setting; the present paper is to our knowledge
the first to make the explicit identification of the non-compact
Lorentzian Dirac as the norm-resolvent limit of a sequence of
truncated Krein spectral triples.

\subsection{Position vs Latr\'emoli\`ere}\label{sec:vs_latremoliere}

Latr\'emoli\`ere's quantum
propinquity~\cite{latremoliere_metric_st_2017,latremoliere2018,
farsi_latremoliere2024,farsi_latremoliere2025} is defined for
\emph{compact} quantum metric spaces. The compact-carrier requirement
enters via the cb-norm finiteness condition (which requires amenable
LCA group structure, satisfied by $\R$) and via the finite Berezin
reach bound (which requires $L^{\infty}$ control on a compact set, NOT
satisfied by $\R$).

A non-compact extension of Latr\'emoli\`ere propinquity would require
addressing the second issue. Possible routes:
\begin{itemize}
\item Locally compact extensions (analogous to how Connes' compact
spectral triples generalize to non-unital triples in the locally
compact setting~\cite{gayral_gracia_bondia_iochum_schucker_varilly2004}):
restrict to $C_{c}^{\infty}$-supported test functions, redefine the
propinquity via approximate identity / proper functions.
\item Functional-calculus version:\ replace $L^{\infty}$-Lipschitz
seminorm with $L^{2}$-quadratic-form seminorm or weighted variants.
\end{itemize}
Either route is multi-month original NCG-math and is not addressed
here.

\subsection{The two-rate hybrid framing in
NCG}\label{sec:hybrid_framing}

The two-rate hybrid convergence in Theorem~\ref{thm:main} —
propinquity rate on the inner arrow, norm-resolvent rate on the outer
arrow — is in the spirit of \emph{multi-scale} or
\emph{two-parameter} convergence frameworks that appear elsewhere in
NCG and mathematical physics:
\begin{itemize}
\item Semiclassical analysis (Sj\"ostrand--Vodev): combined
semiclassical-parameter $\hbar \to 0$ and domain-size parameter
$T \to \infty$ limits.
\item Coarse Bloch decomposition: combined unit-cell-size and
spectral-bandwidth parameters.
\item Continuum limit of lattice gauge theories: combined
lattice-spacing $a \to 0$ and lattice-volume $L \to \infty$ limits.
\end{itemize}
In each case, the two scales must be chosen jointly to give a
well-defined limit. Our admissible scaling condition
(Definition~\ref{def:admissible_scaling}) is the analog in the
Lorentzian-NCG setting.

% =====================================================================
\section{Propinquity-level closure via temporal-Lipschitz-invisibility}\label{sec:g2_metric}
% =====================================================================

\emph{[Descoped (see the Status note): the construction below
quantifies over the degenerate natural-substrate seminorm (whose
kernel far exceeds the scalars; the temporal-Lipschitz-invisibility
identity of Lemma~\ref{lem:extent_zero} is the diagnosis of that
degeneracy, not a feature). The section is retained as a record of
the algebraic construction; its metric-level claims are open pending
the upstream repair.]}

The norm-resolvent outer arrow of Theorem~\ref{thm:main} can be
upgraded to propinquity-level convergence on the non-compact carrier
$\sthree \times \R_{t}$ by combining three ingredients: the
temporal-Lipschitz-invisibility identity of Paper~46
(Lemma~3.2~\cite{paper46}), the Latr\'emoli\`ere pinned proper QMS
hypertopology~\cite{latremoliere2025_hypertopology}, and its
Krein-lift (Paper~48~\cite{paper48}\,\S~3). This \emph{would} close Q1
(G2-metric) on the natural substrate, but the metric level is
\emph{descoped} (it inherits the Paper~45 degeneracy; see the Status
note); only the norm-resolvent / spectral level (\S\S 3--4) is established.

\begin{lemma}[Zero-cost extent element]\label{lem:extent_zero}
Let $\chi_{T} \in C_{c}^{\infty}(\R)$ be a smooth cutoff with
$\chi_{T} \equiv 1$ on $[-(T{-}1)/2,\,(T{-}1)/2]$ and
$\mathrm{supp}\,\chi_{T} \subseteq [-T/2,\,T/2]$.
Define $e_{T} = \mathbf{1}_{S^{3}} \otimes \chi_{T}
\in C_{0}(\sthree \times \R)$. Then
\begin{equation}\label{eq:extent_zero}
   L^{K}(e_{T}) \;=\; \opnorm{[\DL,\, M_{e_{T}}]} \;=\; 0
   \qquad \forall\, T > 0.
\end{equation}
\end{lemma}

\begin{proof}
By the temporal-Lipschitz-invisibility identity (Paper~46
Lemma~3.2~\cite{paper46}), every natural-substrate element
$a = a_{s} \otimes a_{t} \in \Op^{L}$ satisfies
$[\DL, M_{a}] = i[\DGV, M_{a_{s}}] \otimes M_{a_{t}}$.
Setting $a_{s} = \mathbf{1}_{S^{3}}$ (the constant function on $S^{3}$)
and $a_{t} = \chi_{T}$:
$[\DL, M_{e_{T}}] = i[\DGV, M_{\mathbf{1}}] \otimes M_{\chi_{T}} = 0$,
since the Dirac operator commutes with the identity multiplier.
\end{proof}

\begin{remark}\label{rem:extent_mechanism}
Lemma~\ref{lem:extent_zero} is the load-bearing step. In a general
non-compact quantum metric space, the extent element's Lipschitz
norm $L(e) > 0$ restricts the admissible metametric radius via the
condition $L(e) \le r$ (Def.~2.6 of~\cite{latremoliere2025_hypertopology}).
This forces $r$-dependent spatial localization---the source of the
``multi-month'' estimate in the original Q1 formulation. GeoVac's
temporal-Lipschitz-invisibility gives $L^{K}(e_{T}) = 0 \le r$ for
\emph{every} $r > 0$: the temporal extent element is freely available
at zero Lipschitz cost, and the non-compact extension adds nothing
to the propinquity bound.
\end{remark}

\begin{theorem}[G2-metric closure on the natural substrate]\label{thm:g2_metric}
Let $\{(n_{\max}(k),\, N_{t}(k),\, T(k))\}_{k \in \N}$ be an
admissible scaling sequence (Definition~\ref{def:admissible_scaling}).
Then the Krein pinned proper QMS sequence
$\{\Tcal^{L,\mathrm{pp}}_{n_{\max}(k), N_{t}(k), T(k)}\}$ converges
in the Latr\'emoli\`ere hypertopology
(Def.~4.4 of~\cite{latremoliere2025_hypertopology}, Krein-lifted
per Paper~48~\cite{paper48}\,\S~3) to
$\Tcal^{L,\mathrm{pp}}_{S^{3} \times \R_{t}}$:
\begin{equation}\label{eq:g2_metric_bound}
   \eth_{r}^{K}\bigl(\Tcal^{L}_{n_{\max}(k)},\;
   \Tcal^{L}_{S^{3} \times \R}\bigr)
   \;\le\; \Cthreeop(\nmax) \cdot \gamma_{\nmax}
          \;+\; \epsilon(T(k))
   \;\to\; 0
   \qquad \forall\, r > 0,
\end{equation}
where $\Cthreeop(\nmax) = \sqrt{1 - 1/\nmax}$ (Paper~46),
$\gamma_{\nmax} = 2/(\nmax + 1)$ (Paper~45 L2), and
$\epsilon(T) \to 0$ as $T \to \infty$.
\end{theorem}

\begin{proof}
\emph{Step 1 (Continuum limit as Krein pinned proper QMS).}
$\Tcal^{L,\mathrm{pp}}_{S^{3} \times \R}
= (\Acal^{K},\, L^{K},\, \mathcal{M}^{L},\, \omega_{W}^{L})$
with $\Acal^{K} = C_{0}(\sthree \times \R)$,
$L^{K}(f) = \opnorm{[\DL, M_{f}]}$,
$\mathcal{M}^{L} = \mathrm{span}\{M^{\spat}_{N,L,0} \otimes I\}$,
$\omega_{W}^{L}$ the wedge KMS state.
The exhaustive sequence is $h_{n} = \mathbf{1} \otimes \chi_{n}$
with $L^{K}(h_{n}) = 0$
(Lemma~\ref{lem:extent_zero}),
$\omega_{W}^{L}(h_{n}) \to 1$
(thermal tail decay),
$\norm{h_{n}} = 1$.
All nine axioms of~\cite{latremoliere2025_hypertopology}
Defs.~1.18, 1.22, 1.26, 1.29, 1.37, 1.40, 1.42 are verified at
theorem-grade rigor in the Krein-lift
(Paper~48~\cite{paper48}\,\S~3).

\emph{Step 2 (M-tunnel construction).}
For each $k$, construct an M-tunnel
$\tau_{k} = (\mathfrak{D}_{k},\, \pi_{1},\, \pi_{2},\,
L_{\mathfrak{D}},\, e_{k})$ with:
\begin{itemize}
\item $\mathfrak{D}_{k}
      = \Op^{L}_{\nmax(k), \Nt(k)} \oplus C_{0}(\sthree \times \R)$
      (bridging algebra);
\item $\pi_{1} = \mathrm{id}$, \
      $\pi_{2} = B^{\SU(2)}_{\nmax} \otimes (R_{T} \circ B^{\Sone}_{\Nt})$
      (spatial Berezin $\otimes$ temporal restrict-Fej\'er);
\item $L_{\mathfrak{D}}(a_{1}, a_{2})
      = \max(L^{K}_{1}(a_{1}),\, L^{K}_{2}(a_{2}))$;
\item $e_{k} = (I,\, \mathbf{1} \otimes \chi_{T(k)})$
      with $L_{\mathfrak{D}}(e_{k}) = \max(0,\, 0) = 0$
      (Lemma~\ref{lem:extent_zero}).
\end{itemize}

\emph{Step 3 (Reach bound).}
The tunnel reach decomposes into spatial and temporal contributions.
The spatial reach is $\le \Cthreeop \cdot \gamma_{\nmax}$ by
Paper~46 Theorem~5.1 (the compact-carrier propinquity
bound, transported via the M-tunnel restriction to the extent
region $|t| \le T(k)/2$). The temporal contribution is controlled
by the extent mechanism: for states $\phi$ with
$\phi(e_{k}) \ge 1 - \delta$, the mass outside the extent
region $[-T(k)/2,\, T(k)/2]$ is $\le \delta$, giving a
temporal reach $\le \delta$. Setting $\delta = \epsilon(T(k))$
with $\epsilon(T) = 1 - \omega_{W}^{L}(\mathbf{1} \otimes \chi_{T})
\to 0$ (KMS thermal tail):
\[
   \chi(\tau_{k},\, r) \;\le\; \Cthreeop(\nmax) \cdot \gamma_{\nmax}
   \;+\; \epsilon(T(k)).
\]

\emph{Step 4 (Propinquity assembly).}
For each $r > 0$: $L_{\mathfrak{D}}(e_{k}) = 0 \le r$, so
$\tau_{k}$ is admissible at radius $r$. By
Def.~2.23 of~\cite{latremoliere2025_hypertopology} (Krein-lifted):
\[
   \eth_{r}^{K}(\Tcal^{L}_{\nmax(k)},\,
   \Tcal^{L}_{S^{3} \times \R})
   \;\le\; \chi(\tau_{k},\, r)
   \;\le\; \Cthreeop \cdot \gamma_{\nmax} + \epsilon(T(k))
   \;\to\; 0.
\]
Since this holds for every $r > 0$, convergence in the hypertopology
(Def.~4.4) follows.
\end{proof}

\begin{remark}[Bit-exact inheritance from the compact carrier]%
\label{rem:bit_exact_g2}
As $T(k) \to \infty$, the temporal contribution $\epsilon(T(k)) \to 0$
and the bound is dominated by the compact-carrier term
$\Cthreeop \cdot \gamma_{\nmax}$. In particular, the Paper~46 panel
values $\Lprop(2, 3) = 2.0746$, $\Lprop(3, 5) = 1.6101$,
$\Lprop(4, 7) = 1.3223$ are inherited as the asymptotic
non-compact propinquity values. This is the ``free upgrade'' reading
extended from the compact to the non-compact carrier.
\end{remark}

\begin{remark}[Scope: natural substrate only]%
\label{rem:g2_scope}
Theorem~\ref{thm:g2_metric} holds on the \emph{natural substrate}
(chirality-block-diagonal multipliers commuting with $\JL$, Paper~44
Prop.~5.1). On the \emph{enlarged substrate}
(Paper~46 Appendix~B), the temporal-Lipschitz-invisibility identity
fails: $[\DL^{\mathrm{diag}}, M^{\flip} \otimes M^{\temp}] \neq 0$
in general, so $L^{K}(e_{T}) > 0$ and the zero-cost extent mechanism
does not apply. The enlarged-substrate non-compact extension
remains open (Q2 below).
\end{remark}

% =====================================================================
\section{Open questions}\label{sec:open}
% =====================================================================

The construction of \S\S~\ref{sec:inner}--\ref{sec:outer} closes G2
at the norm-resolvent level; the propinquity-level (hypertopology)
upgrade of Theorem~\ref{thm:g2_metric} is descoped (see the Status
note), and the following open questions remain:

\begin{enumerate}
\item[(Q1)] \emph{G2-metric: open (descoped 2026-06-09; see Status note).}
(Theorem~\ref{thm:g2_metric}). The propinquity-level closure uses
temporal-Lipschitz-invisibility (Paper~46 Lemma~3.2) to construct a
zero-cost extent element, reducing the non-compact carrier to the
compact one. The Paper~46 panel values are inherited bit-exactly.

\item[(Q2)] \emph{Enlarged-substrate non-compact extension: open
(descoped 2026-06-09; see Status note).}
The $\JL$-graded gradient norm of Paper~46 Appendix~B gives
$G^{\mathrm{enlarged}}(f) = \norm{\nabla f^{\mathrm{block}}}_{\infty}
+ 2\opnorm{\partial_{t}} \cdot \norm{f^{\mathrm{flip}}}_{\infty}$.
Under admissible scaling, $\opnorm{\partial_{t}} = O(\Nt / T) \to \infty$
(since $T / \Nt \to 0$). The enlarged Lip-1 constraint
$G^{\mathrm{enlarged}} \le 1$ forces
$\norm{f^{\mathrm{flip}}}_{\infty} \le 1/(2\opnorm{\partial_{t}})
\to 0$: the chirality-flip content of the Lip-1 ball is
\emph{squeezed to zero} by the growing temporal Dirac norm.
The enlarged-substrate metametric satisfies
$\eth_{r}^{K,\mathrm{enlarged}} \le \Cthreeop \cdot \gamma_{\nmax}
+ 1/(2\opnorm{\partial_{t}}) + \epsilon(T) \to 0$,
inheriting the Theorem~\ref{thm:g2_metric} bound plus a vanishing
flip-suppression correction. The $O(\Nt/T)$ growth flagged in the
original Q2 formulation is the convergence \emph{mechanism}, not an
obstruction.

\item[(Q3)] \emph{Cross-manifold extension (G3): CLOSED as structural
impossibility.} Paper~32~\cite{paper32}
Theorem~7.4\footnote{Theorem \texttt{thm:g3\_closure} of Paper~32
(\texttt{paper\_32\_spectral\_triple.tex}).}
proves that the cross-manifold tensor product
$\Tcal_{\sthree} \otimes \Tcal^{\mathrm{BL}}_{S^5}$ cannot preserve
both the Riemannian-Dirac character of $S^3$ and the $\pi$-free
Bargmann character of $S^5$ within the standard NCG framework. The
obstruction is categorical (Riemannian vs complex-analytic), tracing
to Bertrand's theorem combined with the Fock and HO rigidity theorems
(Papers~23~\cite{paper23}, 24~\cite{paper24}).

\item[(Q4)] \emph{Inner-factor calibration (G4).} The Yukawa /
Higgs selection question of the Connes--Marcolli almost-commutative
extension is orthogonal to the convergence theorem and remains open.
\end{enumerate}

% =====================================================================
\section*{Acknowledgments}
% =====================================================================

\noindent This work is part of the GeoVac math.OA arc (Papers
38--46). The Lorentzian-extension program was opened by
Sprint~L0~\cite{paper31} (Lorentzian transfer audit) and continued
through Sprints~L1 (Riemannian closure, Paper~42), L2 (Lorentzian
substrate, Paper~43; operator system, Paper~44), L3a-1 (Paper~44),
L3b-2 (Paper~45), L3b-2a--d + L3b-2f-$\beta$ (Paper~46), and L3c-$\alpha$
/ L3c-$\gamma$ (this paper).

% =====================================================================
\begin{thebibliography}{99}
% =====================================================================

\bibitem[Behrndt, Holzmann \& Stelzer-Landauer(2026)]{behrndt_holzmann_stelzer2026}
J.~Behrndt, M.~Holzmann, C.~Stelzer-Landauer,
``Approximation of Dirac operators with $\delta$-shell potentials in the
norm resolvent sense, II:\ Quantitative results,''
\emph{Math. Nachr.} (2026), arXiv:2507.01482.

\bibitem[Bizi, Brouder \& Besnard(2018)]{bizi_brouder_besnard2018}
N.~Bizi, C.~Brouder, and F.~Besnard, ``Space and time dimensions of
algebras with applications to Lorentzian noncommutative geometry and
quantum electrodynamics,'' \emph{J. Math. Phys.} \textbf{59} (2018),
062303. arXiv:1611.07062.

\bibitem[Bunce \& Deddens(1975)]{bunce_deddens1975}
J.~W.~Bunce, J.~A.~Deddens, ``A family of simple C*-algebras related
to weighted shift operators,'' \emph{J. Funct. Anal.} \textbf{19}
(1975), 13--24.

\bibitem[Che, Perales \& Sormani(2025)]{che_perales_sormani2025}
M.~Che, R.~Perales, C.~Sormani,
``Gromov's compactness theorem for the intrinsic timed-Hausdorff distance,''
\emph{Differential Geometry and its Applications} \textbf{103} (2026), arXiv:2510.13069.

\bibitem[Connes(1994)]{connes1994}
A.~Connes, \emph{Noncommutative Geometry}, Academic Press, San Diego,
1994.

\bibitem[Farsi \& Latr\'emoli\`ere(2024)]{farsi_latremoliere2024}
C.~Farsi, F.~Latr\'emoli\`ere, ``The quantum Gromov--Hausdorff
propinquity for crossed products by amenable group actions,''
\emph{J.~Funct.~Anal.} \textbf{286} (2024), 110293.

\bibitem[Farsi \& Latr\'emoli\`ere(2025)]{farsi_latremoliere2025}
C.~Farsi, F.~Latr\'emoli\`ere, ``Continuity for the spectral
propinquity of the Dirac operators associated with an analytic path
of Riemannian metrics,'' arXiv:2504.11715 (April 2025).

\bibitem[Franco \& Eckstein(2014)]{franco_eckstein2014}
N.~Franco, M.~Eckstein, ``Causality in noncommutative geometry,''
\emph{Class. Quant. Grav.} \textbf{30} (2013), 135007 (and follow-ups).

\bibitem[Gayral et al.(2004)]{gayral_gracia_bondia_iochum_schucker_varilly2004}
V.~Gayral, J.~M.~Gracia-Bond\'\i{}a, B.~Iochum, T.~Sch\"ucker,
J.~C.~V\'arilly, ``Moyal planes are spectral triples,''
\emph{Comm. Math. Phys.} \textbf{246} (2004), 569--623.

\bibitem[Geroch(1967)]{geroch1967}
R.~Geroch, ``Topology in general relativity,'' \emph{J. Math. Phys.}
\textbf{8} (1967), 782--786.

\bibitem[Hekkelman \& McDonald(2024)]{hekkelman_mcdonald2024}
E.-M.~Hekkelman, E.~McDonald, ``A noncommutative integral on
spectrally truncated spectral triples, and a link with quantum
ergodicity,'' arXiv:2412.00628 (2024).

\bibitem[Hekkelman \& McDonald(2024b)]{hekkelman_mcdonald2024b}
E.-M.~Hekkelman, E.~McDonald, ``A noncommutative integral on
spectrally truncated spectral triples, and a link with quantum
ergodicity,'' arXiv:2412.00628 (2024).
[Duplicate of \texttt{hekkelman\_mcdonald2024}; an intended second
Hekkelman--McDonald 2024 work could not be located in literature.
PI to consolidate or supply correct second-paper reference.]

\bibitem[Katznelson(2004)]{katznelson2004}
Y.~Katznelson, \emph{An Introduction to Harmonic Analysis}, 3rd ed.,
Cambridge University Press, 2004.

\bibitem[Latr\'emoli\`ere(2018/2023)]{latremoliere_metric_st_2017}
F.~Latr\'emoli\`ere, ``The Gromov--Hausdorff propinquity for
metric spectral triples,'' \emph{Adv. Math.} \textbf{404} (2022),
108393. arXiv:1811.10843.

\bibitem[Latr\'emoli\`ere(2018)]{latremoliere2018}
F.~Latr\'emoli\`ere, ``The quantum Gromov--Hausdorff propinquity,''
\emph{Trans. Amer. Math. Soc.} \textbf{368} (2016), 365--411.
arXiv:1302.4058.

\bibitem[Latr\'emoli\`ere(2025)]{latremoliere2025_hypertopology}
F.~Latr\'emoli\`ere, ``The quantum Gromov--Hausdorff Hypertopology
on the class of pointed Proper Quantum Metric Spaces,''
arXiv:2512.03573 (December 2025).

\bibitem[Mondino \& Sämann(2025)]{mondino_samann2025}
A.~Mondino, C.~Sämann, ``Lorentzian Gromov--Hausdorff
convergence and pre-compactness,'' (2025), arXiv:2504.10380.

\bibitem[Reed \& Simon(1978)]{reed_simon_iv}
M.~Reed, B.~Simon, \emph{Methods of Modern Mathematical Physics
Vol.~IV:\ Analysis of Operators}, Academic Press, New York, 1978.

\bibitem[Strohmaier(2006)]{strohmaier2006}
A.~Strohmaier, ``On noncommutative and pseudo-Riemannian geometry,''
\emph{J. Geom. Phys.} \textbf{56} (2006), 175--195.

% --- Internal GeoVac preprints ---

\bibitem[Loutey, internal Paper~23]{paper23}
J.~Loutey, ``Nuclear shell model qubit Hamiltonians and the Fock
projection rigidity theorem,'' GeoVac internal preprint Paper~23
(2026), at
\texttt{papers/group4\_quantum\_computing/paper\_23\_nuclear\_shell.tex}.

\bibitem[Loutey, internal Paper~24]{paper24}
J.~Loutey, ``The Bargmann--Segal Lattice: $\pi$-Free Discretization
of the Harmonic Oscillator on $S^5$,'' GeoVac internal preprint Paper~24
(2026), at
\texttt{papers/group3\_foundations/paper\_24\_bargmann\_segal.tex}.

\bibitem[Loutey, internal Paper~31]{paper31}
J.~Loutey, ``The Universal/Coulomb Partition of the GeoVac Spectral
Triple,'' GeoVac internal preprint Paper~31 (2026), at
\texttt{papers/group3\_foundations/paper\_31\_universal\_coulomb\_partition.tex}.

\bibitem[Loutey, internal Paper~38]{paper38}
J.~Loutey, ``State-space Gromov--Hausdorff convergence of truncated
Camporesi--Higuchi spectral triples on $S^3$,'' GeoVac internal preprint Paper~38 (2026),
released to Zenodo 2026-05-07. At
\texttt{papers/group1\_operator\_algebras/paper\_38\_su2\_propinquity\_convergence.tex}.

\bibitem[Loutey, internal Paper~39]{paper39}
J.~Loutey, ``Tensor-product propinquity convergence on
$\Tcal_{S^{3}}^{\lambda_{a}}\otimes\Tcal_{S^{3}}^{\lambda_{b}}$,''
GeoVac internal preprint Paper~39 (2026), at
\texttt{papers/group1\_operator\_algebras/paper\_39\_tensor\_propinquity\_convergence.tex}.

\bibitem[Loutey, internal Paper~40]{paper40_unified}
J.~Loutey, ``State-space Gromov--Hausdorff convergence of spectral
truncations of compact Lie groups with bi-invariant metric,'' GeoVac
internal preprint Paper~40 (2026), at
\texttt{papers/group1\_operator\_algebras/paper\_40\_unified\_propinquity\_convergence.tex}.

\bibitem[Loutey, internal Paper~42]{paper42}
J.~Loutey, ``Tomita--Takesaki modular structure on truncated $SU(2)$
spectral triples: four-witness Wick-rotation literal identification
at finite cutoff,'' GeoVac internal preprint Paper~42 (2026), at
\texttt{papers/group1\_operator\_algebras/paper\_42\_modular\_hamiltonian\_four\_witness.tex}.

\bibitem[Loutey, internal Paper~43]{paper43}
J.~Loutey, ``Lorentzian extension of the four-witness Wick-rotation
theorem on truncated $S^3 \times \mathbb{R}$ spectral triples at finite
cutoff,'' GeoVac internal preprint
Paper~43 (2026), at
\texttt{papers/group1\_operator\_algebras/paper\_43\_lorentzian\_extension.tex}.

\bibitem[Loutey, internal Paper~44]{paper44}
J.~Loutey, ``Operator-system extension of the BBB Krein spectral
triple at finite cutoff: propagation number and envelope dependence
on the Camporesi--Higuchi spinor bundle,'' GeoVac internal preprint Paper~44 (2026),
at
\texttt{papers/group1\_operator\_algebras/paper\_44\_lorentzian\_operator\_system.tex}.

\bibitem[Loutey, internal Paper~45]{paper45}
J.~Loutey, ``Lorentzian propinquity on truncated SU(2)
$\times$ U(1)$_{T}$ Krein spectral triples: a degeneracy theorem,'' GeoVac internal
preprint Paper~45 (2026), at
\texttt{papers/group1\_operator\_algebras/paper\_45\_lorentzian\_propinquity.tex}.

\bibitem[Loutey, internal Paper~46]{paper46}
J.~Loutey, ``Strong-form Lorentzian propinquity on
truncated SU(2) $\times$ U(1)$_{T}$ Krein spectral triples,''
GeoVac internal preprint Paper~46 (2026), at
\texttt{papers/group1\_operator\_algebras/paper\_46\_strong\_form\_lorentzian\_propinquity.tex}.

\bibitem[Loutey, internal Paper~32]{paper32}
J.~Loutey, ``The GeoVac Spectral Triple: A Synthesis Paper Locking the
Framework into the Marcolli--van~Suijlekom Gauge-Network Lineage,''
GeoVac internal preprint Paper~32 (2026), at
\texttt{papers/group1\_operator\_algebras/paper\_32\_spectral\_triple.tex}.

\bibitem[Loutey, internal Paper~48]{paper48}
J.~Loutey, ``Krein-pointed proper quantum metric spaces and the bridge
to Mondino--Samann Lorentzian pre-length spaces,''
GeoVac internal preprint Paper~48 (2026), at
\texttt{papers/group1\_operator\_algebras/paper\_48\_krein\_ms\_bridge.tex}.

\end{thebibliography}

\end{document}
